\clearpage
\addcontentsline{toc}{chapter}{KATA PENGANTAR}

% Aturan: Judul 14pt, tebal, kapital. Jarak ke isi 2x1.5 spasi (sekitar 1cm)
\begin{center}
	\textbf{\fontsize{14}{17}\selectfont KATA PENGANTAR}
\end{center}

\vspace{1cm}

Dengan mengucapkan segala puji dan syukur atas kehadirat Allah SWT yang telah memberikan rahmat dan karunia-Nya sehingga saya dapat menyelesaikan skripsi yang berjudul \textbf{"JUDUL SKRIPSI ANDA DITULIS DI SINI"}. 

Adapun tujuan dari skripsi ini adalah sebagai salah satu syarat untuk menyelesaikan Program Sarjana (S1) Jurusan Informatika, Universitas Gunadarma. Dalam kesempatan ini saya ingin menyampaikan ucapan terima kasih atas dukungan dan bantuan yang saya terima, sehingga dapat menyelesaikan penulisan ini, dan perkenankan saya mengucapkan terima kasih kepada:

\begin{enumerate}
	\item Prof. Dr. E. S. Margianti, SE., MM., selaku Rektor Universitas Gunadarma.
	\item Prof. Dr. Ing. Adang Suhendra, SSi., S.Kom., MSc., selaku Dekan Fakultas Teknologi Industri Universitas Gunadarma.
	\item Prof. Dr. Lintang Yuniar Banowosari, S.Kom., MSc., selaku Ketua Program Studi Informatika Universitas Gunadarma.
	\item Bapak Dr. Edi Sukirman, SSi., MM., M.I.Kom., selaku Kepala Bagian Sidang Ujian Universitas Gunadarma.
	\item Nama Pembimbing Lengkap dengan Gelar, selaku dosen pembimbing yang telah memberikan banyak masukan, arahan serta bimbingan kepada saya mulai dari awal sampai selesainya skripsi ini.
	\item Kedua orang tua dan keluarga tercinta yang selalu mendukung dan terus memberikan doa serta motivasi kepada penulis.
	\item Seluruh rekan seperjuangan di Universitas Gunadarma yang telah banyak membantu penulis.
	\item Semua pihak yang tidak dapat disebutkan satu per satu yang telah membantu penyelesaian skripsi ini.
\end{enumerate}

Karena keterbatasan ilmu dan pengalaman yang dimiliki, saya mengharapkan kritik dan saran yang membangun dari pembaca tulisan ini. Saya juga berharap penulisan ini dapat bermanfaat dan menambah wawasan pembaca serta mampu menghasilkan inovasi-inovasi yang baru dari penulisan ini.

Semoga apa yang ada pada penulisan ini dapat bermanfaat bagi kita semua. Aamiin.

\vspace{1.5cm}

% --- BLOK TANDA TANGAN ---
\begin{flushright}
	Jakarta, 27 Februari 2026 \\ % Sesuaikan tanggalnya
	\vspace{0.2cm}
	
	% CARA MEMASUKKAN TTD:
	% 1. Siapkan file gambar ttd Anda (misal: ttd.png), pastikan backgroundnya transparan.
	% 2. Taruh di folder yang sama dengan file .tex ini.
	% 3. Hapus tanda persen (%) di baris bawah ini untuk mengaktifkan gambarnya:
	% \includegraphics[height=2cm]{ttd.png} 
	
	\vspace{1.5cm} % Hapus atau kecilkan nilai vspace ini jika gambar ttd sudah diaktifkan agar jaraknya tidak terlalu jauh
	
	( Nama Lengkap Anda )
\end{flushright}